% \iffalse meta-comment
%
% Copyright (C) 2024 by Lukas Heindl <oss.heindl+latex@protonmail.com>
% ---------------------------------------------------------------------------
% This work may be distributed and/or modified under the
% conditions of the LaTeX Project Public License, either version 1.3
% of this license or (at your option) any later version.
% The latest version of this license is in
%   http://www.latex-project.org/lppl.txt
% and version 1.3 or later is part of all distributions of LaTeX
% version 2005/12/01 or later.
%
% This work has the LPPL maintenance status `maintained'.
%
% The Current Maintainer of this work is Lukas Heindl.
%
% This work consists of the files beamercolorthemeCatppuccin.dtx and beamercolorthemeCatppuccin.ins
% and the derived filebase beamercolorthemeCatppuccin.sty.
%
% \fi
%
% \iffalse
%<*driver>
\ProvidesFile{timberborndrawing.dtx}
%</driver>
%<package>\NeedsTeXFormat{LaTeX2e}[1999/12/01]
%<package>\ProvidesPackage{timberborndrawing}
%<*package>
   [2026/03/16 v0.1.0 initial draft]
%</package>
%
%<*driver>
\documentclass{ltxdoc}
\usepackage[a4paper,margin=25mm,left=50mm,nohead]{geometry}
\usepackage{babel}

\usepackage{tcolorbox}
\tcbuselibrary{skins,listingsutf8,documentation}

\usepackage{catppuccinpalette}
\usepackage{listings}
\lstset{
	rulecolor=\color{black},
	emphstyle=\color{CtpYellow},
	basicstyle=\tiny\ttfamily,
	stringstyle=\color{CtpGreen},
	commentstyle=\color{CtpOverlay2},
	keywordstyle=\color{CtpMauve},
	backgroundcolor=\color{CtpBase},
	frame=none,
	%
	columns=fullflexible, % columns of chars
	tabsize=2,
	%
	captionpos=b,
	%
	breakautoindent=true, %
	breakindent=1em,
	% postbreak=\mbox{\textcolor{red}{$\hookrightarrow$}\space},
	breaklines=true,
	%
	showstringspaces=false,
	showtabs=false,
	showspaces=false,
	keepspaces=true,
	%
	inputencoding=utf8, extendedchars=true,
	literate=
	{á}{{\'a}}1 {é}{{\'e}}1 {í}{{\'i}}1 {ó}{{\'o}}1 {ú}{{\'u}}1
	{Á}{{\'A}}1 {É}{{\'E}}1 {Í}{{\'I}}1 {Ó}{{\'O}}1 {Ú}{{\'U}}1
	{à}{{\`a}}1 {è}{{\`e}}1 {ì}{{\`i}}1 {ò}{{\`o}}1 {ù}{{\`u}}1
	{À}{{\`A}}1 {È}{{\'E}}1 {Ì}{{\`I}}1 {Ò}{{\`O}}1 {Ù}{{\`U}}1
	{ä}{{\"a}}1 {ë}{{\"e}}1 {ï}{{\"i}}1 {ö}{{\"o}}1 {ü}{{\"u}}1
	{Ä}{{\"A}}1 {Ë}{{\"E}}1 {Ï}{{\"I}}1 {Ö}{{\"O}}1 {Ü}{{\"U}}1
	{â}{{\^a}}1 {ê}{{\^e}}1 {î}{{\^i}}1 {ô}{{\^o}}1 {û}{{\^u}}1
	{Â}{{\^A}}1 {Ê}{{\^E}}1 {Î}{{\^I}}1 {Ô}{{\^O}}1 {Û}{{\^U}}1
	{Ã}{{\~A}}1 {ã}{{\~a}}1 {Õ}{{\~O}}1 {õ}{{\~o}}1
	{œ}{{\oe}}1 {Œ}{{\OE}}1 {æ}{{\ae}}1 {Æ}{{\AE}}1 {ß}{{\ss}}1
	{ű}{{\H{u}}}1 {Ű}{{\H{U}}}1 {ő}{{\H{o}}}1 {Ő}{{\H{O}}}1
	{ç}{{\c c}}1 {Ç}{{\c C}}1 {ø}{{\o}}1 {å}{{\r a}}1 {Å}{{\r A}}1
	{€}{{\euro}}1 {£}{{\pounds}}1 {«}{{\guillemotleft}}1
	{»}{{\guillemotright}}1 {ñ}{{\~n}}1 {Ñ}{{\~N}}1 {¿}{{?`}}1
}

\usepackage{tikz}
\usetikzlibrary{graphdrawing,shapes}
\usetikzlibrary{arrows.meta,automata,arrows}
\usegdlibrary{layered}

\usepackage{adjustbox}

\usepackage{timberborndrawing}[2026/03/16]

\RecordChanges
\begin{document}
  \DocInput{\jobname.dtx}
  \PrintChanges
\end{document}
%</driver>
% \fi
%
% \CheckSum{362}
%
% \CharacterTable
%  {Upper-case    \A\B\C\D\E\F\G\H\I\J\K\L\M\N\O\P\Q\R\S\T\U\V\W\X\Y\Z
%   Lower-case    \a\b\c\d\e\f\g\h\i\j\k\l\m\n\o\p\q\r\s\t\u\v\w\x\y\z
%   Digits        \0\1\2\3\4\5\6\7\8\9
%   Exclamation   \!     Double quote  \"     Hash (number) \#
%   Dollar        \$     Percent       \%     Ampersand     \&
%   Acute accent  \'     Left paren    \(     Right paren   \)
%   Asterisk      \*     Plus          \+     Comma         \,
%   Minus         \-     Point         \.     Solidus       \/
%   Colon         \:     Semicolon     \;     Less than     \<
%   Equals        \=     Greater than  \>     Question mark \?
%   Commercial at \@     Left bracket  \[     Backslash     \\
%   Right bracket \]     Circumflex    \^     Underscore    \_
%   Grave accent  \`     Left brace    \{     Vertical bar  \|
%   Right brace   \}     Tilde         \~}
%
%
% \changes{v0.0.0}{2026/03/16}{First draft}
%
% \providecommand*{\url}{\texttt}
% \GetFileInfo{timberborndrawing.dtx}
% \title{\textsf{timberborndrawing} --- Draw graph for timberborn
% }
% \author{Lukas Heindl \\ \url{oss.heindl+latex@protonmail.com} \\Github:
% \url{https://github.com/atticus-sullivan/timberborn_drawing/}}
% \date{\fileversion~from \filedate}
%
% \maketitle
%
% \begin{abstract}
% The package provides a small interface between \LaTeX{}
% and a Lua backend that reads JSON data and generates
% TikZ drawings representing Timberborn production chains.
% 
% Documents using this package need to be compiled with
% LuaLaTeX. The package requires |tikz| (and |pgfkeys|).
% \end{abstract}
%
% \tableofcontents
%
% \section{Usage}
% Load this package as |\usepackage{timberborndrawing}|
%
% The overall principle of this is the following:
% First, all information is read in by the lua layer (usually from a json file).
% Then, this information can be drawn.
% For this the lua layer emits the corresponding tikz code which needs to be set within a \verb|tikzpicture|.
% The emitted code only defines the content (nodes and connections).
% The appearance can be modified e.g. by adjusting the tikz styles.
%
% Note that the graph is only in a way that connects a \emph{hut} with consumed and produced \emph{resource}s.
% Resources are never connected directly with each other. Huts are also never connected directly with each other.
%
% The following commands are defined:
% \begin{docCommand}[]{timberborndrawingReset}{}
%   Reset the internal Lua state.
% \end{docCommand}
%
% \begin{docCommand}[]{timberborndrawingLoadJSON}{\marg{file-path}}
%   Read the data from the specified JSON file.
%   Usually this JSON file is generated automatically \footnote{\url{https://github.com/atticus-sullivan/timberborn_plots}}
% \end{docCommand}
%
% \begin{docCommand}[]{timberborndrawingDrawAll}{\marg{showCnt}}
%   Emit the tikz code.
%   Must be called inside a \verb|tikzpicture|.
% \end{docCommand}
%
% \subsection{Keys used for adjusting}
% The following keys can be used to customize the output.
%
% \tcbmakedocSubKey{docTimberborndrawing}{timberborndrawing}
%
% \begin{docTimberborndrawing}{hut/node}{\meta{style}}{style, initially empty}
%   This style is applied to the nodes representing a \emph{hut} (something which has an input and an output).
% \end{docTimberborndrawing}
%
% \begin{docTimberborndrawing}{resource/node}{\meta{style}}{style, initially empty}
%    This style is applied to the nodes representing a \emph{resource} (something which can be stored).
% \end{docTimberborndrawing}
% \begin{docTimberborndrawing}{resource/pos}{\meta{style}}{style, initially empty}
%    When a resource is net-positive (currently builds up).
% \end{docTimberborndrawing}
% \begin{docTimberborndrawing}{resource/neg}{\meta{style}}{style, initially empty}
%    When a resource is net-negative (currently empties).
% \end{docTimberborndrawing}
%
% \begin{docTimberborndrawing}{edgeIn}{\meta{style}}{style, initially empty}
%    This style is applied to the incoming edges of a hut.
% \end{docTimberborndrawing}
% \begin{docTimberborndrawing}{edgeOut}{\meta{style}}{style, initially empty}
%    This style is applied to the outgoing edges of a hut.
% \end{docTimberborndrawing}
% \begin{docTimberborndrawing}{every edge}{\meta{style}}{style, initially empty}
%    This style is applied to all edges in the graph.
% \end{docTimberborndrawing}
%
% \subsection{Example}
%<*doc>
\begin{filecontents}[noheader]{data.json}
[
  {
    "template": "Birch",
    "name": "Birch",
    "recipe": null,
    "cnt": 420,
    "dur": "7d",
    "inputs": {},
    "outputs": {
      "Wood": 1
    }
  },
  {
    "template": "LumberMill.Folktails",
    "name": "LumberMill.Folktails",
    "recipe": null,
    "cnt": 2,
    "dur": "1.3h",
    "inputs": {
      "Wood": 1,
      "Power": 50
    },
    "outputs": {
      "Planks": 1
    }
  }
]
\end{filecontents}
\begin{dispListing*}{listing options={language=json},title={\verb|data.json|}}
[
  {
    "template": "Birch",
    "name": "Birch",
    "recipe": null,
    "cnt": 42,
    "dur": "7d",
    "inputs": {},
    "outputs": {
      "Wood": 1
    }
  },
  {
    "template": "LumberMill.Folktails",
    "name": "LumberMill.Folktails",
    "recipe": null,
    "cnt": 2,
    "dur": "1.3h",
    "inputs": {
      "Wood": 1,
      "Power": 50
    },
    "outputs": {
      "Planks": 1
    }
  }
]
\end{dispListing*}

\begin{dispExample*}{listing options={language=[LaTeX]TeX}}
\timberborndrawingReset
\timberborndrawingLoadJSON{./data.json}
\maxsizebox{\linewidth}{\textheight}{
  \begin{tikzpicture}[
    /timberborndrawing/hut/node/.style={
      draw,ellipse,
      align=center,
    },
    /timberborndrawing/resource/neg/.style={CtpRed},
    /timberborndrawing/resource/pos/.style={CtpGreen},
    /timberborndrawing/every edge/.style={
      draw,
      ->,>={Stealth[scale=2]},shorten >=1pt,
    },
    layered layout,
    level distance=15ex,
    sibling sep=8em,
    ]
    \timberborndrawingDrawAll{true}
  \end{tikzpicture}
}
\end{dispExample*}
%</doc>
% 
%
% \newpage
% \section{Implementation}
% This package uses |timberborndrawing| as prefix/directory where
% possible. Since this is not possible for latex macro names, in this
% occasions |timberborndrawing@| is used as prefix.
% \subsection{timberborndrawing.sty}
% \subsubsection{Global Stuff}
%    \begin{macrocode}
%<*package>
\NeedsTeXFormat{LaTeX2e}
\ProvidesPackage{timberborndrawing}[2026/03/16 Timberborn drawing helper]

\RequirePackage{tikz}
\usetikzlibrary{graphs,graphdrawing}
\RequirePackage{amsmath}
\RequirePackage{pgfkeys}
%    \end{macrocode}
%
% \subsubsection{Key Definitions}
%
% The pgfkeys interface used by the Lua backend.
%
%    \begin{macrocode}

\pgfkeys{
	/timberborndrawing/hut/node/.style={},
	/timberborndrawing/resource/node/.style={},
	/timberborndrawing/resource/pos/.style={},
	/timberborndrawing/resource/neg/.style={},
	/timberborndrawing/every edge/.style={},
	/timberborndrawing/edgeIn/.style={},
	/timberborndrawing/edgeOut/.style={},
}
%    \end{macrocode}
% \subsubsection{Lua Backend}
%
% Load the Lua module which contains the drawing logic.
%
%    \begin{macrocode}
\directlua{
tim = require("timberborndrawing")
}
%    \end{macrocode}
%
% \subsubsection{User Commands}
%
% Reset the Lua drawing state.
%
%    \begin{macrocode}
\newcommand{\timberborndrawingReset}{
  \directlua{tim.reset()}
}
%    \end{macrocode}
%
% Load a JSON analysis file.
%
%    \begin{macrocode}
\newcommand{\timberborndrawingLoadJSON}[1]{
  \directlua{tim.loadJSON("#1")}
}
%    \end{macrocode}
%
% Execute the drawing routine.
%
%    \begin{macrocode}
\NewDocumentCommand{\timberborndrawingDrawAll}{m}{
  \directlua{tim.drawAll(#1)}
}
%</package>
%    \end{macrocode}
\endinput
%
% \iffalse
%</package>
% \fi
%
% \Finale
\endinput
